\section{Week 4: Maximum Likelihood Estimation}

\begin{quote}
\textbf{Chapter 2: Parameter estimation methods} · Lecture 1 h · Lab 2 h
\textbf{Running data:} the Week 2--3 forgetting experiment (binomial
recall counts), plus a new SDT experiment with three payoff conditions.
\end{quote}

\begin{center}\rule{0.5\linewidth}{0.5pt}\end{center}

\subsection{Key terms}

Important terms of this lesson, with their Persian equivalents (repeated
even if they appeared in earlier weeks).

{\def\LTcaptype{none} % do not increment counter
\begin{longtable}[]{@{}
  >{\raggedright\arraybackslash}p{(\linewidth - 4\tabcolsep) * \real{0.3333}}
  >{\raggedright\arraybackslash}p{(\linewidth - 4\tabcolsep) * \real{0.3333}}
  >{\raggedright\arraybackslash}p{(\linewidth - 4\tabcolsep) * \real{0.3333}}@{}}
\toprule\noalign{}
\begin{minipage}[b]{\linewidth}\raggedright
English
\end{minipage} & \begin{minipage}[b]{\linewidth}\raggedright
\RL{\CMfa فارسی}
\end{minipage} & \begin{minipage}[b]{\linewidth}\raggedright
\RL{\CMfa توضیح کوتاه}
\end{minipage} \\
\midrule\noalign{}
\endhead
\bottomrule\noalign{}
\endlastfoot
likelihood & \RL{\CMfa درست‌نمایی} & \RL{\CMfa احتمال داده به‌عنوان تابعی از
پارامترها} \\
estimate / estimator & \RL{\CMfa برآورد} / \RL{\CMfa برآوردگر} & \RL{\CMfa مقدار به‌دست‌آمده
برای پارامتر} / \RL{\CMfa روش به دست آوردن آن} \\
parameter & \RL{\CMfa پارامتر} & \RL{\CMfa کمیت قابل تنظیم مدل که بین افراد یا شرایط
فرق می‌کند} \\
forgetting curve & \RL{\CMfa منحنی فراموشی} & \RL{\CMfa کاهش یادآوری با گذشت
زمان} \\
binomial distribution & \RL{\CMfa توزیع دوجمله‌ای} & \RL{\CMfa توزیع تعداد موفقیت در}
n \RL{\CMfa کوشش} \\
signal detection theory (SDT) & \RL{\CMfa نظریه‌ی تشخیص علامت} & \RL{\CMfa مدلی که
حساسیت را از سوگیری پاسخ جدا می‌کند} \\
log-likelihood / NLL & \RL{\CMfa لگاریتم درست‌نمایی} / \RL{\CMfa منفی آن} &
\RL{\CMfa مقداری که بهینه‌سازها کمینه یا بیشینه می‌کنند} \\
least squares & \RL{\CMfa کمترین مربعات} & \RL{\CMfa روش برازش با کمینه کردن مجموع
مربع خطاها} \\
maximum likelihood estimate (MLE) & \RL{\CMfa برآورد بیشینه‌ی درست‌نمایی} &
\RL{\CMfa پارامترهایی که داده را محتمل‌ترین می‌کنند} \\
noise & \RL{\CMfa نویز} & \RL{\CMfa تغییرات تصادفی در شواهد یا در داده} \\
fit / fitting & \RL{\CMfa برازش} & \RL{\CMfa تنظیم پارامترها تا مدل با داده جور
شود} \\
deviance (G²) & \RL{\CMfa انحراف} (deviance) & \RL{\CMfa معیار عدم برازش نسبت به مدل
اشباع‌شده} \\
bounds / constraints & \RL{\CMfa کران‌ها} / \RL{\CMfa قیدها} & \RL{\CMfa محدودیت‌های مجاز
برای مقدار پارامترها} \\
response bias & \RL{\CMfa سوگیری پاسخ} & \RL{\CMfa تمایل به یک پاسخ، مستقل از
محرک} \\
sensitivity (d′) & \RL{\CMfa حساسیت} & \RL{\CMfa فاصله‌ی توزیع علامت و نویز} \\
trial & \RL{\CMfa کوشش} (trial) & \RL{\CMfa یک بار انجام تکلیف توسط شرکت‌کننده} \\
optimizer & \RL{\CMfa بهینه‌ساز} & \RL{\CMfa الگوریتمی که کمینه‌ی تابع را پیدا
می‌کند} \\
SSE / RMSD & \RL{\CMfa مجموع مربعات خطا} / \RL{\CMfa ریشه‌ی میانگین مربع انحراف} &
\RL{\CMfa معیارهای فاصله‌ی داده و پیش‌بینی} \\
saturated model & \RL{\CMfa مدل اشباع‌شده} & \RL{\CMfa مدلی که داده را دقیقاً بازتولید
می‌کند} \\
likelihood-ratio test (LRT) & \RL{\CMfa آزمون نسبت درست‌نمایی} & \RL{\CMfa مقایسه‌ی دو
مدل تودرتو} \\
exponential / power function & \RL{\CMfa تابع نمایی} / \RL{\CMfa توانی} & \RL{\CMfa دو شکل
رایج برای منحنی فراموشی و یادگیری} \\
liberal / conservative & \RL{\CMfa آسان‌گیر} / \RL{\CMfa محافظه‌کار} & \RL{\CMfa تمایل به
گفتن «بله»} / \RL{\CMfa «نه»} \\
token & token \RL{\CMfa (واحد متن)} & \RL{\CMfa کلمه یا بخشی از کلمه که مدل زبانی
پیش‌بینی می‌کند} \\
perplexity & perplexity \RL{\CMfa (سردرگمی)} & e \RL{\CMfa به توان} NLL \RL{\CMfa هر}
token \\
\end{longtable}
}

\begin{center}\rule{0.5\linewidth}{0.5pt}\end{center}

\subsection{Learning outcomes}

By the end of this week, students will be able to:

\begin{enumerate}
\def\labelenumi{\arabic{enumi}.}
\tightlist
\item
  \textbf{Distinguish} probability (data varying, parameters fixed) from
  likelihood (data fixed, parameters varying).
\item
  \textbf{Derive} the likelihood and log-likelihood for binomial
  (Bernoulli) and Gaussian data models, and \textbf{show} that least
  squares is maximum likelihood under equal-variance Gaussian noise.
\item
  \textbf{Implement} maximum likelihood estimation (MLE) by minimizing
  the negative log-likelihood (NLL) in Python, for forgetting curves and
  for a multi-condition SDT model.
\item
  \textbf{Evaluate} absolute goodness of fit with the deviance \(G^2\)
  and \textbf{compare} MLE with least-squares estimates.
\item
  \textbf{Interpret} constrained models (e.g., ``payoffs change bias,
  not sensitivity'') as substantive psychological hypotheses.
\end{enumerate}

\begin{center}\rule{0.5\linewidth}{0.5pt}\end{center}

\subsection{Motivation \& big picture}

In Week 3, least squares treated every data point as equally reliable.
But a recall proportion of 0.50 measured on 50 items is \emph{noisier}
than a proportion of 0.95: binomial variance \(f(1-f)/n\) depends on
\(f\). Least squares also has nothing to say about data that are not
continuous numbers, such as a single yes/no response, a choice among
options, or a confidence rating.

Maximum likelihood replaces ``make predictions close to data'' with a
sharper question:

\begin{quote}
\textbf{Which parameter values make the data we actually observed most
probable?}
\end{quote}

To answer it, the model must say how probable \emph{any} possible
dataset is, which is exactly the noise model (component 4 of Week 2's
anatomy). Once we have it, MLE handles counts, choices, RTs and mixtures
of these in the same way. It also gives us, for free, an absolute
measure of misfit (the deviance), the basis for standard errors (Week
5), information criteria (Week 7), and the likelihood term in Bayes'
rule (Weeks 9--13). \textbf{The likelihood is the single most important
object in this course.}

\paragraph{Where is this used?}

\begin{itemize}
\tightlist
\item
  \textbf{Almost every published cognitive model} is fitted by maximum
  likelihood or by its Bayesian extension (Weeks 9--13). When a paper
  reports ``model fits'', ``AIC'' or ``log-likelihood'', MLE is
  underneath.
\item
  \textbf{Choice data}: which option a person picks (in a memory test, a
  gamble, a learning task) is not a continuous number, so least squares
  does not apply. The likelihood handles choices naturally.
\item
  \textbf{Psychophysics}: the psychometric function (probability of
  detecting a stimulus as a function of its strength) is fitted by
  binomial MLE in every vision and hearing lab.
\item
  \textbf{Testing hypotheses inside a model}: ``does the payoff change
  sensitivity or only bias?'' becomes a comparison of two likelihoods
  (Section 6).
\item
  \textbf{AI}: large language models are trained by maximizing the
  likelihood of text (see the AI \& LLMs section).
\end{itemize}

\begin{center}\rule{0.5\linewidth}{0.5pt}\end{center}

\subsection{Lecture notes}

\paragraph{In plain words}

Imagine trying different parameter values one by one. For each value,
ask the model: ``if this were the true value, how probable would the
data I actually collected be?'' That number is the \textbf{likelihood}
of that parameter value. The \textbf{maximum likelihood estimate (MLE)}
is the parameter value under which the observed data would have been
most probable. Example: someone recalls 7 of 10 words. If their true
recall probability were 0.5, getting exactly 7 would be fairly unlikely
(0.117); if it were 0.7, it would be more likely (0.267). So 0.7 is the
better explanation, and indeed the MLE is 7/10. In practice we work with
the \textbf{logarithm} of the likelihood (so products become sums) and
we let the computer \emph{minimize} its negative (NLL), using the
optimizers of Week 3. Two bonuses: least squares turns out to be a
special case of MLE, and the likelihood gives an absolute test of misfit
(\(G^2\)).

\textbf{Suggested timing (60 min):} probability vs.~likelihood 10 ·
binomial MLE 15 · log-likelihood and NLL 10 · Gaussian case (LS = MLE)
10 · goodness of fit \(G^2\) 10 · constrained models 5.

\subsubsection{1. Probability vs.~likelihood}

The same function \(p(\mathbf{y}\mid\boldsymbol\theta)\) can be read in
two ways:

{\def\LTcaptype{none} % do not increment counter
\begin{longtable}[]{@{}
  >{\raggedright\arraybackslash}p{(\linewidth - 8\tabcolsep) * \real{0.2000}}
  >{\raggedright\arraybackslash}p{(\linewidth - 8\tabcolsep) * \real{0.2000}}
  >{\raggedright\arraybackslash}p{(\linewidth - 8\tabcolsep) * \real{0.2000}}
  >{\raggedright\arraybackslash}p{(\linewidth - 8\tabcolsep) * \real{0.2000}}
  >{\raggedright\arraybackslash}p{(\linewidth - 8\tabcolsep) * \real{0.2000}}@{}}
\toprule\noalign{}
\begin{minipage}[b]{\linewidth}\raggedright
Reading
\end{minipage} & \begin{minipage}[b]{\linewidth}\raggedright
Fixed
\end{minipage} & \begin{minipage}[b]{\linewidth}\raggedright
Varies
\end{minipage} & \begin{minipage}[b]{\linewidth}\raggedright
Name
\end{minipage} & \begin{minipage}[b]{\linewidth}\raggedright
Sums/integrates to 1?
\end{minipage} \\
\midrule\noalign{}
\endhead
\bottomrule\noalign{}
\endlastfoot
Before the experiment & \(\boldsymbol\theta\) & \(\mathbf{y}\) &
Probability (mass/density) of data & Yes, over \(\mathbf{y}\) \\
After the experiment & \(\mathbf{y}\) (observed) & \(\boldsymbol\theta\)
& \textbf{Likelihood}
\(L(\boldsymbol\theta) = p(\mathbf{y}\mid\boldsymbol\theta,\mathcal{M})\)
& \textbf{No}, not over \(\boldsymbol\theta\) \\
\end{longtable}
}

The likelihood is \textbf{not} the probability that the parameters are
correct. It is the probability of the data \emph{under} each candidate
parameter value. (Turning it into a probability distribution over
\(\boldsymbol\theta\) requires a prior: Bayes, Week 9.)

\textbf{Example.} A participant recalls \(y = 7\) of \(n = 10\) items.
Under a binomial model with recall probability \(\theta\):

\[
L(\theta) = \binom{10}{7}\theta^{7}(1-\theta)^{3}.
\]

\(L(0.5) = 0.117\), \(L(0.7) = 0.267\): the observed data are about 2.3
times more probable if \(\theta = 0.7\) than if \(\theta = 0.5\).

\paragraph{Reading the math}

\(L(\theta) = \binom{10}{7}\theta^{7}(1-\theta)^{3}\):

\begin{itemize}
\tightlist
\item
  \(\theta^7\): the probability that 7 specific items are recalled (each
  with probability \(\theta\)).
\item
  \((1-\theta)^3\): the probability that the other 3 are forgotten.
\item
  \(\binom{10}{7} = 120\) (``10 choose 7''): the number of different
  ways to pick which 7 of the 10 items are the recalled ones. It does
  not depend on \(\theta\), so it never changes \emph{where} the maximum
  is.
\item
  Plug in \(\theta = 0.7\):
  \(120 \times 0.7^7 \times 0.3^3 = 120 \times 0.0824 \times 0.027 = 0.267\).
\end{itemize}

The key mental switch: in probability, \(\theta\) is fixed and we ask
about different data; in likelihood, the data are fixed (we already have
them) and we compare different \(\theta\).

\subsubsection{2. The maximum likelihood estimate}

\[
\hat{\boldsymbol\theta}_{\text{ML}} = \arg\max_{\boldsymbol\theta} L(\boldsymbol\theta) = \arg\max_{\boldsymbol\theta} \ell(\boldsymbol\theta),
\qquad \ell(\boldsymbol\theta) = \ln L(\boldsymbol\theta).
\]

\textbf{Why the log?} 1. With independent observations the likelihood is
a \textbf{product}, \(L = \prod_i p(y_i\mid\boldsymbol\theta)\), which
underflows to 0 on a computer for even a few hundred trials. The log
turns it into a \textbf{sum},
\(\ell = \sum_i \ln p(y_i\mid\boldsymbol\theta)\). 2. The log is
strictly increasing, so it has the same maximizer. 3. Sums are easier to
differentiate.

Optimizers \emph{minimize}, so in code we minimize the \textbf{negative
log-likelihood},
\(\mathrm{NLL}(\boldsymbol\theta) = -\ell(\boldsymbol\theta)\).

\textbf{Analytic example (binomial).}
\(\ell(\theta) = \text{const} + y\ln\theta + (n-y)\ln(1-\theta)\).
Setting the derivative to zero:

\[
\frac{d\ell}{d\theta} = \frac{y}{\theta} - \frac{n-y}{1-\theta} = 0 \quad\Rightarrow\quad \hat\theta = \frac{y}{n}.
\]

The MLE of a probability is the observed proportion, which is
reassuring. For cognitive models the predicted probability is a
\emph{function} of the parameters,
\(\theta_i = f(x_i;\boldsymbol\theta)\), and we usually need a numerical
optimizer.

\subsubsection{3. The binomial likelihood for forgetting curves}

With \(y_i\) items recalled out of \(n_i\) at interval \(t_i\), and
independent intervals:

\[
\ell(\boldsymbol\theta) = \sum_{i=1}^{N}\Big[\ln\binom{n_i}{y_i} + y_i \ln f(t_i;\boldsymbol\theta) + (n_i - y_i)\ln\big(1 - f(t_i;\boldsymbol\theta)\big)\Big].
\]

\begin{itemize}
\tightlist
\item
  The binomial coefficient does not depend on \(\boldsymbol\theta\), so
  it does not affect the MLE (but keep it if you want to compare
  log-likelihoods across different data codings;
  \nolinkurl{scipy.stats.binom.logpmf} includes it).
\item
  Observations predicted to be near 0 or 1 have small variance, so a
  prediction that misses them is punished more heavily. MLE performs the
  variance weighting of Week 3's weighted least squares
  \emph{automatically}.
\item
  \textbf{Numerical safety:} if the model predicts \(f = 0\) where
  \(y > 0\), then \(\ln 0 = -\infty\). Clip predictions to
  \([\epsilon, 1-\epsilon]\) with, e.g., \(\epsilon = 10^{-10}\).
\end{itemize}

\textbf{Bernoulli vs.~binomial.} If we keep the individual 0/1
responses, \(\ell = \sum_{\text{trials}} [y\ln f + (1-y)\ln(1-f)]\).
This differs from the binomial version only by the constant binomial
coefficients, so it gives the same MLE.

\paragraph{Reading the math}

The forgetting log-likelihood is the single-proportion log-likelihood
from Section 2, written once per retention interval and added up:

\begin{itemize}
\tightlist
\item
  for interval \(i\): \(y_i \ln f(t_i)\) rewards predicting a high
  recall probability where many items were recalled;
  \((n_i - y_i)\ln(1 - f(t_i))\) rewards predicting a low recall
  probability where many were forgotten.
\item
  Adding over intervals is allowed because the intervals are assumed
  \textbf{independent}: the log of a product of probabilities is the sum
  of the logs.
\item
  Why it weights points automatically: if the model predicts
  \(f = 0.95\) but only 40 of 50 are recalled, the term
  \(10 \ln(0.05) \approx -30\) is very costly. Near \(f = 0.5\), the
  same absolute miss costs much less. Points that \emph{should} be
  precise are trusted more.
\end{itemize}

\subsubsection{4. The Gaussian case: least squares is maximum
likelihood}

Suppose \(y_i = f(x_i;\boldsymbol\theta) + \varepsilon_i\) with
\(\varepsilon_i \sim \mathcal{N}(0, \sigma)\), independent and with the
\textbf{same} \(\sigma\). Then

\[
\ell(\boldsymbol\theta, \sigma) = \sum_{i=1}^{N}\ln\left[\frac{1}{\sqrt{2\pi}\sigma}\exp\left(-\frac{(y_i - f(x_i;\boldsymbol\theta))^2}{2\sigma^2}\right)\right]
= -\frac{N}{2}\ln(2\pi\sigma^2) - \frac{\mathrm{SSE}(\boldsymbol\theta)}{2\sigma^2}.
\]

For any fixed \(\sigma\), maximizing \(\ell\) over \(\boldsymbol\theta\)
means minimizing SSE. So \textbf{least squares = MLE under i.i.d.
equal-variance Gaussian noise}. Maximizing over \(\sigma\) as well gives
\(\hat\sigma^2 = \mathrm{SSE}/N\).

This explains \emph{when} least squares is appropriate (continuous,
roughly symmetric noise with constant variance) and when it is not
(counts, proportions, choices, skewed RTs).

\begin{quote}
\textbf{Advanced: properties of the MLE.} Under regularity conditions
(identifiable model, true parameters not on the boundary, \(N\) large),
the MLE is \emph{consistent} (converges to the true value),
\emph{asymptotically normal}, and \emph{asymptotically efficient} (no
consistent estimator has smaller variance in the limit). Its sampling
variance is approximated by the inverse of the \textbf{Fisher
information}, the curvature of \(-\ell\) at its minimum. Week 5 uses
exactly this to compute standard errors. In small samples the MLE can be
biased (e.g., \(\hat\sigma^2 = \mathrm{SSE}/N\) underestimates
\(\sigma^2\)).
\end{quote}

\subsubsection{\texorpdfstring{5. Absolute goodness of fit: the deviance
\(G^2\)}{5. Absolute goodness of fit: the deviance G\^{}2}}

The log-likelihood value alone has no absolute meaning. To judge it,
compare the model with the \textbf{saturated model}, which has one free
probability per data cell and therefore reproduces every observed
proportion exactly (\(\hat f_i = y_i/n_i\)). The \textbf{deviance} is

\[
G^2 = 2\big[\ell_{\text{saturated}} - \ell(\hat{\boldsymbol\theta})\big]
= 2\sum_i \left[y_i \ln\frac{y_i}{n_i \hat f_i} + (n_i - y_i)\ln\frac{n_i - y_i}{n_i(1-\hat f_i)}\right],
\]

with \(0 \ln 0 = 0\). If the model is correct and counts are not too
small, \(G^2\) approximately follows a \(\chi^2\) distribution with
\(\text{df} = (\text{number of cells}) - k\). A large \(G^2\) (small
\(p\)) means the model misfits \emph{more than sampling noise can
explain}.

Caveats: with very large \(N\), any small misfit becomes
``significant''; with small \(N\), real misfit can go undetected. A
non-significant \(G^2\) does \textbf{not} prove the model is right (Week
6).

\paragraph{Reading the math}

\begin{itemize}
\tightlist
\item
  The \textbf{saturated model} is the ``cheating'' model that predicts
  every observed proportion exactly. Its log-likelihood is the best any
  model could possibly achieve on these data.
\item
  \(G^2 = 2(\ell_{\text{saturated}} - \ell(\hat{\boldsymbol\theta}))\)
  measures how far our model falls short of that ceiling. \(G^2 = 0\)
  would be a perfect fit.
\item
  \textbf{Degrees of freedom} = number of data cells − number of fitted
  parameters. With 8 intervals and 2 parameters, df = 6.
\item
  If the model is true, \(G^2\) behaves like a \(\chi^2\) variable with
  that df, whose average is df. So \(G^2 \approx 6\) with df = 6 is a
  normal amount of misfit (just noise); \(G^2 = 26\) with df = 6
  (\(p = .0002\)) means the model systematically misses the data.
\end{itemize}

\subsubsection{6. Constrained models as hypotheses}

MLE makes it easy to fit \textbf{several conditions at once with some
parameters shared}. Example: an SDT experiment with three payoff
conditions (rewarding ``yes'' responses, neutral, rewarding ``no'').

{\def\LTcaptype{none} % do not increment counter
\begin{longtable}[]{@{}
  >{\raggedright\arraybackslash}p{(\linewidth - 6\tabcolsep) * \real{0.2500}}
  >{\raggedright\arraybackslash}p{(\linewidth - 6\tabcolsep) * \real{0.2500}}
  >{\raggedright\arraybackslash}p{(\linewidth - 6\tabcolsep) * \real{0.2500}}
  >{\raggedright\arraybackslash}p{(\linewidth - 6\tabcolsep) * \real{0.2500}}@{}}
\toprule\noalign{}
\begin{minipage}[b]{\linewidth}\raggedright
Model
\end{minipage} & \begin{minipage}[b]{\linewidth}\raggedright
Free parameters
\end{minipage} & \begin{minipage}[b]{\linewidth}\raggedright
\(k\)
\end{minipage} & \begin{minipage}[b]{\linewidth}\raggedright
Hypothesis
\end{minipage} \\
\midrule\noalign{}
\endhead
\bottomrule\noalign{}
\endlastfoot
Shared sensitivity & \(d'\); \(c_1, c_2, c_3\) & 4 & Payoffs change
\emph{bias} only \\
Separate & \(d'_1, d'_2, d'_3\); \(c_1, c_2, c_3\) & 6 & Payoffs may
also change \emph{sensitivity} (this is the saturated model for 6
rates) \\
\end{longtable}
}

The joint log-likelihood is the sum over conditions of the binomial
log-likelihoods of hits and false alarms. Comparing the two models'
maximized log-likelihoods tests a substantive question. With nested
models like these, \(2(\ell_{6} - \ell_{4})\) is approximately
\(\chi^2\) with 2 df: the \textbf{likelihood-ratio test}, developed in
Week 7.

\begin{quote}
\textbf{Advanced: SDT MLE and the formulas.} For a single condition, the
MLE of \((d', c)\) from counts is exactly what the Week 2 formulas give
\emph{without} correction, because the model is saturated (2 parameters,
2 rates) and the MLE reproduces the observed rates. The log-linear
correction is \emph{not} MLE; it is closer to a Bayesian estimate with a
weak prior (Week 9).
\end{quote}

\paragraph{For an engineer}

\begin{itemize}
\tightlist
\item
  This is \textbf{maximum-likelihood estimation} exactly as in
  estimation theory (ML channel or frequency estimation, ML sequence
  detection). The only change is the noise model: binomial counts
  instead of AWGN.
\item
  \textbf{LS = ML under AWGN} (Section 4) is a result you already use:
  with white Gaussian noise of known, equal variance, the ML estimator
  minimizes the squared error.
\item
  The \textbf{Fisher information} and the asymptotic variance of the MLE
  (Advanced box) are the same quantities as in the \textbf{Cramér--Rao
  lower bound}. Week 5 uses them to compute standard errors.
\item
  Comparing a constrained and a free model with
  \(2(\ell_{\text{free}} - \ell_{\text{constrained}})\) is the
  \textbf{generalized likelihood-ratio test (GLRT)}: both hypotheses
  have unknown parameters that are replaced by their ML estimates.
\item
  \textbf{Payoffs and criterion shifts}: in Bayesian detection the
  optimal threshold moves with the costs and priors. Section 6 tests
  whether \emph{people} move only their threshold (\(c\)) when the
  payoffs change, while their ``SNR'' (\(d'\)) stays the same.
\end{itemize}

\paragraph{For the instructor}

\begin{itemize}
\tightlist
\item
  \textbf{The one sentence to repeat:} ``the likelihood is the
  probability of the data we observed, as a function of the
  parameters''. Write it on the board and leave it there all session.
\item
  \textbf{Plot it.} Draw \(L(\theta)\) for 7 of 10 (worked example, Part
  1) as a curve over \(\theta\) from 0 to 1. Ask: where is the peak?
  Does the area under the curve equal 1? (No.) Then draw it for 70 of
  100: same peak, much narrower curve. That narrowness is precision, the
  topic of Week 5.
\item
  \textbf{Use your background.} Many students find MLE abstract; you can
  present it as ``the estimation method engineers use when the noise is
  not Gaussian''.
\item
  \textbf{Typical questions.}

  \begin{itemize}
  \tightlist
  \item
    \emph{``Is the likelihood the probability that \(\theta\) is
    correct?''} No.~That needs a prior (Week 9).
  \item
    \emph{``Why minimize NLL instead of maximizing \(L\)?''} Same
    answer, but sums of logs avoid numerical underflow, and optimizers
    are written to minimize.
  \item
    \emph{``Is a non-significant \(G^2\) proof the model is right?''}
    No: it only means the misfit is not larger than the noise, and with
    few data, real misfit can hide.
  \end{itemize}
\item
  \textbf{Time management.} Section 4 (Gaussian case) can be shown
  quickly as ``LS is a special case''. Spend the saved time on Section 5
  (\(G^2\)).
\end{itemize}

\begin{center}\rule{0.5\linewidth}{0.5pt}\end{center}

\subsection{Worked example}

\subsubsection{Part 1: likelihood of a single proportion}

\(y = 7\) recalled out of \(n = 10\):

{\def\LTcaptype{none} % do not increment counter
\begin{longtable}[]{@{}
  >{\raggedright\arraybackslash}p{(\linewidth - 8\tabcolsep) * \real{0.2000}}
  >{\raggedright\arraybackslash}p{(\linewidth - 8\tabcolsep) * \real{0.2000}}
  >{\raggedright\arraybackslash}p{(\linewidth - 8\tabcolsep) * \real{0.2000}}
  >{\raggedright\arraybackslash}p{(\linewidth - 8\tabcolsep) * \real{0.2000}}
  >{\raggedright\arraybackslash}p{(\linewidth - 8\tabcolsep) * \real{0.2000}}@{}}
\toprule\noalign{}
\begin{minipage}[b]{\linewidth}\raggedright
\(\theta\)
\end{minipage} & \begin{minipage}[b]{\linewidth}\raggedright
0.5
\end{minipage} & \begin{minipage}[b]{\linewidth}\raggedright
0.6
\end{minipage} & \begin{minipage}[b]{\linewidth}\raggedright
\textbf{0.7}
\end{minipage} & \begin{minipage}[b]{\linewidth}\raggedright
0.8
\end{minipage} \\
\midrule\noalign{}
\endhead
\bottomrule\noalign{}
\endlastfoot
\(L(\theta) = \binom{10}{7}\theta^7(1-\theta)^3\) & 0.117 & 0.215 &
\textbf{0.267} & 0.201 \\
\(\ell(\theta)\) & −2.144 & −1.537 & \textbf{−1.321} & −1.603 \\
\end{longtable}
}

The maximum is at \(\hat\theta = 7/10 = 0.7\). Note that the likelihood
values do not sum to 1 across \(\theta\): likelihood is not a
probability distribution over parameters.

\subsubsection{Part 2: a two-interval forgetting experiment}

\(y = (30, 13)\) items recalled out of \(n = 50\) at \(t = (1, 7)\)
days. Model: \(f(t) = A e^{-\lambda t}\).

\textbf{Step 1: NLL for a candidate.} For
\((A, \lambda) = (0.70, 0.15)\): \(f(1) = 0.7e^{-0.15} = 0.6025\),
\(f(7) = 0.7e^{-1.05} = 0.2450\).

\[
\ell = \ln\mathrm{Binom}(30; 50, 0.6025) + \ln\mathrm{Binom}(13; 50, 0.2450) = -2.167 + (-2.088) = -4.255,
\]

so NLL \(= 4.255\).

\textbf{Step 2: compare.} For \((0.70, 0.12)\): \(f = (0.621, 0.302)\)
and NLL \(= 4.488\), which is worse. The first candidate makes the data
more probable.

\textbf{Step 3: the MLE.} With 2 parameters and 2 data points the model
can match both observed proportions exactly: \(A e^{-\lambda} = 0.60\)
and \(A e^{-7\lambda} = 0.26\). Dividing gives
\(e^{6\lambda} = 0.60/0.26\), so \(\hat\lambda = \ln(2.308)/6 = 0.139\)
and \(\hat A = 0.60\,e^{0.139} = 0.690\), with NLL \(= 4.225\) (the
minimum). A numerical optimizer gives the same values.

\textbf{Interpretation.} The MLE says this participant recalls about
69\% at the start and loses about 13\% of what remains per day
(\(1 - e^{-0.139}\)). Because \(k = N\) here, \(G^2 = 0\): the model
\emph{cannot} misfit, so the perfect fit carries no evidence for the
exponential form. Only data with more cells than parameters can test a
model.

\begin{center}\rule{0.5\linewidth}{0.5pt}\end{center}

\subsection{Lab (2 hours)}

\subsubsection{Goals}

\begin{enumerate}
\def\labelenumi{\arabic{enumi}.}
\tightlist
\item
  Plot a likelihood function and see why we work with log-likelihoods.
\item
  Fit exponential and power forgetting models by MLE; compare with least
  squares; evaluate with \(G^2\).
\item
  Verify numerically that Gaussian MLE equals least squares.
\item
  Fit a constrained multi-condition SDT model and \textbf{interpret}
  what it says about payoffs.
\end{enumerate}

\textbf{Timing:} Part A 15 min · Part B 35 min · Part C 15 min · Part D
35 min · Part E 20 min (exercises).

\subsubsection{Setup}

\begin{quote}
\textbf{Code · \RL{\CMfa کد:}} \nolinkurl{code/week04.py}, cell 1 · \RL{\CMfa سلول}
1
\end{quote}

\subsubsection{Part A: the likelihood function (15 min)}

\begin{quote}
\textbf{Code · \RL{\CMfa کد:}} \nolinkurl{code/week04.py}, cell 2 · \RL{\CMfa سلول}
2
\end{quote}

\textbf{Expected output:}

\Needspace{6\baselineskip}
\begin{Verbatim}
Area under L(theta): 0.091 (not 1: L is not a distribution over theta)
Product of 2000 trial probabilities: 0.0
Sum of their logs: -1239.52
\end{Verbatim}

\subsubsection{Part B: MLE for forgetting curves (35 min)}

\begin{quote}
\textbf{Code · \RL{\CMfa کد:}} \nolinkurl{code/week04.py}, cell 3 · \RL{\CMfa سلول}
3
\end{quote}

\textbf{Expected output:}

\Needspace{9\baselineskip}
\begin{Verbatim}
model method      A   rate     NLL    SSE
  exp    MLE 0.6377 0.0458 29.8123 0.1100
  exp     LS 0.6568 0.0566 30.6983 0.1061
  pow    MLE 0.8773 0.4214 19.8303 0.0262
  pow     LS 0.8712 0.4251 19.8585 0.0259
True (power): A = 0.9, d = 0.5
\end{Verbatim}

\textbf{What to notice:} - Each method wins on its \emph{own} criterion:
MLE has the lowest NLL, LS the lowest SSE. - For the (true) power model
the two methods agree closely. - For the exponential model the two
methods differ more (decay rate 0.046 vs.~0.057). The binomial
likelihood gives more weight to the precisely measured points near the
floor (long delays), so the ML curve decays more slowly in order to stay
closer to them.

\begin{quote}
\textbf{Warning: a real optimizer trap.} Replace \nolinkurl{"L-BFGS-B"}
with \nolinkurl{"Nelder-Mead"} in the ML fit above (keeping
\nolinkurl{bounds=bounds}) and rerun. For the exponential model, scipy's
bounded Nelder--Mead collapses against the bound \nolinkurl{A = 1} and
reports \nolinkurl{success} with NLL = 79.5, far worse than the true
minimum of 29.8 (it is even worse than the LS estimate's NLL). Bounded
simplex methods can degenerate when vertices are pushed onto a bound.
Two quick checks would catch this: an estimate sitting exactly on a
bound, and a ``MLE'' whose NLL is \emph{higher} than that of some other
parameter values you have tried. Multi-start (Week 3) is the general
cure.
\end{quote}

Now evaluate absolute fit with the deviance:

\begin{quote}
\textbf{Code · \RL{\CMfa کد:}} \nolinkurl{code/week04.py}, cell 4 · \RL{\CMfa سلول}
4
\end{quote}

\textbf{Expected output:}

\Needspace{5\baselineskip}
\begin{Verbatim}
exp: G2 =  26.27, df = 6, p = 0.0002
pow: G2 =   6.31, df = 6, p = 0.3895
\end{Verbatim}

\textbf{Interpret:} the exponential model misfits far more than binomial
noise allows (\(p = .0002\)): we can \textbf{reject} it as a description
of these data. The power model's misfit is consistent with sampling
noise. That does not prove it is the right model; it means these data
give no reason to reject it.

\begin{quote}
\textbf{Code · \RL{\CMfa کد:}} \nolinkurl{code/week04.py}, cell 5 · \RL{\CMfa سلول}
5
\end{quote}

\subsubsection{Part C: Gaussian MLE equals least squares (15 min)}

Treat the proportions as continuous and assume Gaussian noise with a
free \(\sigma\) (a third parameter). The estimates of \(A\) and \(d\)
should match least squares, and \(\hat\sigma\) should equal
\(\sqrt{\mathrm{SSE}/N}\).

\begin{quote}
\textbf{Code · \RL{\CMfa کد:}} \nolinkurl{code/week04.py}, cell 6 · \RL{\CMfa سلول}
6
\end{quote}

\textbf{Expected output:}

\Needspace{5\baselineskip}
\begin{Verbatim}
Gaussian MLE: A = 0.8712, d = 0.4251, sigma = 0.0569
Least squares: A = 0.8711, d = 0.4251, sqrt(SSE/N) = 0.0569
\end{Verbatim}

\subsubsection{Part D: a constrained multi-condition SDT model (35 min)}

\textbf{The experiment.} Participants detect a faint tone under three
payoff schemes: \textbf{liberal} (hits rewarded), \textbf{neutral}, and
\textbf{conservative} (false alarms penalized). 80 signal and 80 noise
trials per condition. We simulate one participant whose sensitivity is
the same in all conditions and whose criterion shifts with the payoffs.

\begin{quote}
\textbf{Code · \RL{\CMfa کد:}} \nolinkurl{code/week04.py}, cell 7 · \RL{\CMfa سلول}
7
\end{quote}

\textbf{Expected output:}

\Needspace{12\baselineskip}
\begin{Verbatim}
      condition  hits  n_signal  fas  n_noise
0       liberal    73        80   33       80
1       neutral    58        80   16       80
2  conservative    44        80    7       80

Shared d':   d' = 1.496, c = [-0.555  0.124  0.618], NLL = 13.134
Separate d': d' = [1.577 1.439 1.482], c = [-0.568  0.122  0.615], NLL = 13.043

Likelihood-ratio statistic = 0.18, df = 2, p = 0.913
\end{Verbatim}

Note that the estimated \nolinkurl{d'} (about 1.5) is higher than the
generating value of 1.2. With 80 trials per cell that is ordinary
sampling error. Week 5 shows how to put a confidence interval around it.

\textbf{Interpret (write in your notebook):} 1. Does constraining \(d'\)
to be equal across payoff conditions cost much likelihood? What do you
conclude about the effect of payoffs? 2. The separate-\(d'\) estimates
differ from each other by a few tenths. Why is that not evidence that
payoffs change sensitivity? 3. Check that the separate-\(d'\) model
reproduces the observed hit and false-alarm rates exactly. Why must it?

\subsubsection{Exercises (increasing difficulty)}

\textbf{Exercise 1 \ensuremath{\star}.} Derive by hand the MLE of \(\theta\) for a
Bernoulli sequence with \(y\) successes in \(n\) trials, and check it
numerically with \nolinkurl{minimize_scalar} for \(y = 13\), \(n = 40\).

\textbf{Exercise 2 \ensuremath{\star}.} In Part B, refit both models with the
\textbf{Bernoulli} likelihood on item-level data (expand each count into
0/1 responses). Show that the estimates are identical to the binomial
fits and that the NLLs differ by a constant. What is the constant?

\textbf{Exercise 3 \ensuremath{\star}\ensuremath{\star} (zero counts).} Set the last two counts to 0
(\nolinkurl{recalled[-2:] = 0}). Which fitting methods from Weeks 2--4
still work? Fit both models by MLE and comment on what happens to the
estimates and to \(G^2\).

\textbf{Exercise 4 \ensuremath{\star}\ensuremath{\star} (a third SDT model).} Fit a model in which the
criterion is the \emph{same} in all conditions but \(d'\) differs
(\(k = 4\)). Compare its NLL with the shared-\(d'\) model. Both have
\(k = 4\); why can you compare them directly by NLL here? Which one is
psychologically more plausible for a payoff manipulation?

\textbf{Exercise 5 \ensuremath{\star}\ensuremath{\star}\ensuremath{\star} (Advanced: LS vs.~MLE by simulation).} Simulate
500 datasets from the power model (\(A = 0.9\), \(d = 0.5\), the 8
intervals, \(n = 50\)), fit each by LS and by MLE, and compare the two
estimators of \(d\): mean (bias), SD, and root-mean-squared error. Which
is better here? Repeat with \(n = 10\). (This is good preparation for
A2.)

\subsubsection{Lab check-off}

Show: (1) the LS-vs-MLE table and \(G^2\) results from Part B, (2) the
Gaussian-vs-LS comparison from Part C, and (3) your answers to the
interpretation questions of Part D.

\begin{center}\rule{0.5\linewidth}{0.5pt}\end{center}

\subsection{Connection to AI \& LLMs}

\textbf{LLMs are trained by maximum likelihood.} The ``cross-entropy
loss'' used to train every language model is the negative log-likelihood
of the next token, averaged over a huge text corpus. Training an LLM is
therefore an MLE problem with billions of parameters, solved by gradient
descent instead of Nelder--Mead. Two quantities you meet in AI papers
are this week's quantities in disguise: \textbf{loss} = NLL per token,
and \textbf{perplexity} = \(e^{\text{NLL per token}}\), the effective
number of equally likely choices the model is ``confused'' between.

The simplest language model, a \textbf{bigram model}, has a closed-form
MLE: \(P(w_2\mid w_1)\) = count(\(w_1 w_2\))/count(\(w_1\)). It shows a
familiar problem. Any word pair never seen in training gets probability
0, so a single new pair makes the likelihood of a whole sentence 0 (NLL
= \ensuremath{\infty}). Smoothing solves this, in the same way that the log-linear
correction of Week 2 solved hit rates of 1.

\begin{quote}
\textbf{Code · \RL{\CMfa کد:}} \nolinkurl{code/week04.py}, cell 8 · \RL{\CMfa سلول}
8
\end{quote}

\textbf{Expected output:}

\Needspace{8\baselineskip}
\begin{Verbatim}
P(cat | the) = 0.333
'the dog sat on the mat .'   MLE      NLL =  3.584  perplexity =   1.82
'the dog sat on the mat .'   add-one  NLL =  9.537  perplexity =   4.90
'the dog saw the cat .'      MLE      NLL =    inf  perplexity =    inf
'the dog saw the cat .'      add-one  NLL =  9.624  perplexity =   6.85
\end{Verbatim}

\textbf{Try it (no API needed).} - \ensuremath{\star} Which word pair in ``the dog saw
the cat .'' was never seen in the corpus? Check by hand that the MLE of
\(P(\text{cat}\mid\text{the})\) is 1/3. - \ensuremath{\star}\ensuremath{\star} Smoothing makes the
\emph{training} sentence less likely (NLL 3.6 \ensuremath{\rightarrow} 9.5) but makes the new
sentence possible. Replace add-one with add-\(k\) and choose \(k\) by
leave-one-sentence-out cross-validation (a preview of Week 6). - \ensuremath{\star}\ensuremath{\star}
Discussion: the MLE of a bigram model reproduces the training corpus as
closely as possible. Why is ``maximizing likelihood on internet text''
not the same as ``learning to tell the truth''?

\begin{center}\rule{0.5\linewidth}{0.5pt}\end{center}

\subsection{Common mistakes \& misconceptions}

\begin{enumerate}
\def\labelenumi{\arabic{enumi}.}
\tightlist
\item
  \textbf{``The likelihood is the probability that the parameters are
  true.''} It is the probability (density) of the \emph{data} given the
  parameters. It does not integrate to 1 over \(\boldsymbol\theta\).
\item
  \textbf{Multiplying probabilities instead of summing logs.} The
  product underflows to 0 for realistic numbers of trials (Part A).
\item
  \textbf{Maximizing when the optimizer minimizes.} \nolinkurl{minimize}
  needs the \textbf{negative} log-likelihood. Forgetting the minus sign
  sends the optimizer to the \emph{worst} parameters.
\item
  \textbf{Not clipping probabilities.} Predicted probabilities of
  exactly 0 or 1 produce \(-\infty\) log-likelihoods and crash the
  optimizer.
\item
  \textbf{Comparing NLLs across different datasets or data codings.} NLL
  depends on the number of observations and on constants (e.g., binomial
  coefficients). Compare only models fitted to the same data, coded the
  same way.
\item
  \textbf{Using least squares for choices or counts by default.} It
  ignores the data's noise structure. Use the likelihood that matches
  how the data were generated.
\item
  \textbf{Reading a non-significant \(G^2\) as proof the model is
  correct.} It only says the misfit is not detectably larger than
  sampling noise at this sample size.
\item
  \textbf{Forgetting that a saturated or exactly-determined model always
  fits perfectly} (worked example, Part D separate-\(d'\) model). A
  perfect fit is then uninformative.
\end{enumerate}

\begin{center}\rule{0.5\linewidth}{0.5pt}\end{center}

\subsection{Check your understanding}

\textbf{1.} In one sentence each: what varies and what is fixed in a
probability distribution vs.~a likelihood function? \textbf{Answer}
Probability: parameters fixed, data vary, and it sums/integrates to 1
over data. Likelihood: data fixed at the observed values, parameters
vary, and it need not integrate to 1 over parameters.

\textbf{2.} Give two reasons for working with log-likelihoods.
\textbf{Answer} (i) Numerical: products of many probabilities underflow,
while sums of logs do not. (ii) The log is monotone, so the maximizer is
unchanged, and sums are easier to differentiate and optimize.

\textbf{3.} Under what assumptions is least squares equivalent to
maximum likelihood? Show the key step. \textbf{Answer} Independent
Gaussian errors with equal variance \(\sigma^2\). Then
\(\ell = -\tfrac{N}{2}\ln(2\pi\sigma^2) - \mathrm{SSE}/(2\sigma^2)\), so
for any \(\sigma\), maximizing \(\ell\) over \(\boldsymbol\theta\) is
minimizing SSE.

\textbf{4.} A model with 2 parameters is fitted to 8 binomial cells and
gives \(G^2 = 4.1\). What are the df, and what do you conclude?
\textbf{Answer} df = 8 − 2 = 6. A \(\chi^2_6\) has mean 6, so 4.1 is
well within the range expected from sampling noise (\(p \approx 0.66\)).
There is no evidence of misfit, but that does not prove the model
correct.

\textbf{5.} Why does MLE with a binomial likelihood weight data points
differently than plain least squares does? \textbf{Answer} The binomial
log-likelihood reflects the variance \(f(1-f)/n\): points predicted near
0 or 1 (small variance) are fitted more tightly, and points with larger
\(n\) count more. Plain least squares weights every squared error
equally.

\textbf{6.} In the payoff experiment, what psychological hypothesis does
the shared-\(d'\) model express, and how would you test it?
\textbf{Answer} That payoffs change only response bias (criterion), not
perceptual sensitivity. Test it by comparing the maximized
log-likelihood with the separate-\(d'\) model:
\(2(\ell_6 - \ell_4) \sim \chi^2_2\) approximately (likelihood-ratio
test), or with AIC/BIC (Week 7).

\textbf{7.} You fit a model and the optimizer returns NLL = −35.2
(negative). Is something wrong? \textbf{Answer} Not necessarily. For
continuous data, densities can exceed 1 (e.g., a narrow Gaussian), so
log-densities can be positive and the NLL negative. For discrete data
(binomial), probabilities are \ensuremath{\leq} 1, so the NLL must be \ensuremath{\geq} 0, and a
negative value would indicate a bug.

\begin{center}\rule{0.5\linewidth}{0.5pt}\end{center}

\subsection{Assignment}

\subsubsection{A2: Least squares vs.~maximum likelihood (1.0 point of
20; released Week 4, due at the start of the Week 6 lecture)}

\textbf{Data.} Generate the dataset with exactly this code (everyone
gets the same data):

\begin{quote}
\textbf{Code · \RL{\CMfa کد:}} \nolinkurl{code/week04.py}, cell 9 · \RL{\CMfa سلول}
9
\end{quote}

\textbf{Tasks.}

\begin{enumerate}
\def\labelenumi{\arabic{enumi}.}
\tightlist
\item
  \textbf{Fit (2 pts).} Fit the exponential and power models to these
  data by (a) least squares on proportions and (b) binomial MLE. Report
  all estimates in one table with SSE and NLL.
\item
  \textbf{Evaluate (2 pts).} Compute \(G^2\), df and \(p\) for each
  model's MLE fit. Plot the data with 95\% binomial intervals, both
  fitted curves, and residuals. Which model, if any, can be rejected?
\item
  \textbf{Simulation study (3 pts).} Treat the power model with
  \(A = 0.95\), \(d = 0.7\) as the truth. Simulate 500 datasets with the
  same design, fit each with LS and MLE, and report the bias, SD and
  RMSE of \(\hat d\) for each method. Show histograms. Repeat with
  \(n = 100\) items and describe what changes.
\item
  \textbf{Interpret (2 pts, max. 300 words).} Which estimator would you
  recommend for recall data, and why? Relate your answer to the noise
  structure of binomial data and to the near-floor counts (1--2 of 20)
  at the longest delays. What would happen to the Week 2 linearization
  method if one of those counts were 0?
\item
  \textbf{Code quality (1 pt).} Reproducible, commented notebook; seeds
  fixed; functions reused rather than copy-pasted.
\end{enumerate}

Late submissions lose 10\% per day (max. 3 days).

\begin{center}\rule{0.5\linewidth}{0.5pt}\end{center}

\subsection{Further reading}

\begin{itemize}
\tightlist
\item
  \textbf{Farrell \& Lewandowsky (2018), chapter 4.} Maximum likelihood
  parameter estimation: likelihood functions for common data types,
  log-likelihood, and fitting in practice. \emph{(Main reading.)}
\item
  \textbf{Azimifar \& Ahmadi (1396)}, \emph{Shenasayi-ye amari-ye olgu},
  vol.~1 (in Persian): the section on maximum likelihood parameter
  estimation, with the statistical-learning view {[}verify chapter{]}.
\item
  \textbf{Lee \& Wagenmakers (2014), chapter 3.} Inferences with
  binomials. Read it now for the binomial set-up; we revisit it from the
  Bayesian side in Week 9.
\end{itemize}

\emph{Preview for Week 5:} an estimate without an uncertainty is half a
result. We compute standard errors and confidence intervals from the
curvature of the likelihood and by bootstrapping, and we discover how
averaging participants can create a ``law'' that no individual follows.
